% Options for packages loaded elsewhere
\PassOptionsToPackage{unicode}{hyperref}
\PassOptionsToPackage{hyphens}{url}
\documentclass[
  spanish,
  12pt,
  letterpaper]{article}
\usepackage{xcolor}
\usepackage[margin=1.55cm]{geometry}
\usepackage{amsmath,amssymb}
\setcounter{secnumdepth}{-\maxdimen} % remove section numbering
\usepackage{iftex}
\ifPDFTeX
  \usepackage[T1]{fontenc}
  \usepackage[utf8]{inputenc}
  \usepackage{textcomp} % provide euro and other symbols
\else % if luatex or xetex
  \usepackage{unicode-math} % this also loads fontspec
  \defaultfontfeatures{Scale=MatchLowercase}
  \defaultfontfeatures[\rmfamily]{Ligatures=TeX,Scale=1}
\fi
\usepackage{lmodern}
\ifPDFTeX\else
  % xetex/luatex font selection
  \setmainfont[]{Andika}
\fi
% Use upquote if available, for straight quotes in verbatim environments
\IfFileExists{upquote.sty}{\usepackage{upquote}}{}
\IfFileExists{microtype.sty}{% use microtype if available
  \usepackage[]{microtype}
  \UseMicrotypeSet[protrusion]{basicmath} % disable protrusion for tt fonts
}{}
\makeatletter
\@ifundefined{KOMAClassName}{% if non-KOMA class
  \IfFileExists{parskip.sty}{%
    \usepackage{parskip}
  }{% else
    \setlength{\parindent}{0pt}
    \setlength{\parskip}{6pt plus 2pt minus 1pt}}
}{% if KOMA class
  \KOMAoptions{parskip=half}}
\makeatother
\usepackage{graphicx}
\makeatletter
\newsavebox\pandoc@box
\newcommand*\pandocbounded[1]{% scales image to fit in text height/width
  \sbox\pandoc@box{#1}%
  \Gscale@div\@tempa{\textheight}{\dimexpr\ht\pandoc@box+\dp\pandoc@box\relax}%
  \Gscale@div\@tempb{\linewidth}{\wd\pandoc@box}%
  \ifdim\@tempb\p@<\@tempa\p@\let\@tempa\@tempb\fi% select the smaller of both
  \ifdim\@tempa\p@<\p@\scalebox{\@tempa}{\usebox\pandoc@box}%
  \else\usebox{\pandoc@box}%
  \fi%
}
% Set default figure placement to htbp
\def\fps@figure{htbp}
\makeatother
\ifLuaTeX
\usepackage[bidi=basic,shorthands=off]{babel}
\else
\usepackage[bidi=default,shorthands=off]{babel}
\fi
\ifPDFTeX
\else
\babelfont{rm}[]{Andika}
\fi
\ifLuaTeX
  \usepackage{selnolig} % disable illegal ligatures
\fi
\setlength{\emergencystretch}{3em} % prevent overfull lines
\providecommand{\tightlist}{%
  \setlength{\itemsep}{0pt}\setlength{\parskip}{0pt}}
\usepackage{multicol}
\usepackage{enumitem}
\usepackage{xcolor}
\definecolor{azul}{HTML}{476885}
\pagestyle{empty}
\setlength{\parindent}{0pt}
\setlength{\parskip}{3pt}
\setlist[enumerate]{leftmargin=*,itemsep=4pt,topsep=3pt}
\newcommand{\titulo}[1]{\vspace{5pt}{\large\bfseries\color{azul}#1}\par\vspace{3pt}}
\usepackage{bookmark}
\IfFileExists{xurl.sty}{\usepackage{xurl}}{} % add URL line breaks if available
\urlstyle{same}
\hypersetup{
  pdftitle={Taller: uso de tablas estadísticas},
  pdflang={es},
  hidelinks,
  pdfcreator={LaTeX via pandoc}}

\title{Taller: uso de tablas estadísticas}
\usepackage{etoolbox}
\makeatletter
\providecommand{\subtitle}[1]{% add subtitle to \maketitle
  \apptocmd{\@title}{\par {\large #1 \par}}{}{}
}
\makeatother
\subtitle{Teoría de Probabilidad \textbar{} Normal, t de Student,
chi-cuadrado y F}
\author{}
\date{\vspace{-2.5em}}

\begin{document}
\maketitle

\vspace{-1.1cm}

\textbf{Nombre:}

\rule{9cm}{0.4pt}

\hfill **Fecha:**

\rule{2.4cm}{0.4pt}

\vspace{5pt}{\large\bfseries\color{azul}Cara A. Lectura de tablas y cálculo de percentiles}\par\vspace{3pt}

\textbf{Instrucciones.} Use exclusivamente las tablas estadísticas
suministradas. En todos los ejercicios, encuentre el valor crítico
indicado y señale la probabilidad que consultó. Use la notación \(q_p\)
tal que \(P(X\leq q_p)=p\). Si una tabla presenta probabilidades de cola
derecha, convierta \(p\) en \(1-p\).

\vspace{5pt}{\large\bfseries\color{azul}1. Distribución normal estándar: $Z\sim N(0,1)$}\par\vspace{3pt}

\begin{multicols}{2}
\begin{enumerate}[label=\alph*)]
\item $q_{0.90}$
\item $q_{0.95}$
\item $q_{0.975}$
\item $q_{0.99}$
\item $q_{0.10}$
\item $q_{0.025}$
\end{enumerate}
\end{multicols}

\vspace{5pt}{\large\bfseries\color{azul}2. Distribución normal: $X\sim N(120,20^2)$}\par\vspace{3pt}

Calcule los percentiles usando \(x_p=120+20z_p\).

\begin{multicols}{2}
\begin{enumerate}[label=\alph*)]
\item $x_{0.05}$
\item $x_{0.25}$
\item $x_{0.90}$
\item $x_{0.975}$
\end{enumerate}
\end{multicols}

\vspace{5pt}{\large\bfseries\color{azul}3. Distribución t de Student}\par\vspace{3pt}

Encuentre \(t_{\nu,p}\), donde \(\nu\) son los grados de libertad.

\begin{multicols}{2}
\begin{enumerate}[label=\alph*)]
\item $\nu=8,\ p=0.95$
\item $\nu=12,\ p=0.975$
\item $\nu=20,\ p=0.99$
\item $\nu=30,\ p=0.90$
\item $\nu=15,\ p=0.025$
\item $\nu=10,\ p=0.05$
\end{enumerate}
\end{multicols}

\vspace{5pt}{\large\bfseries\color{azul}4. Distribución chi-cuadrado}\par\vspace{3pt}

Encuentre \(\chi^2_{\nu,p}\); distinga los percentiles inferior y
superior.

\begin{multicols}{2}
\begin{enumerate}[label=\alph*)]
\item $\nu=5,\ p=0.95$
\item $\nu=10,\ p=0.975$
\item $\nu=15,\ p=0.05$
\item $\nu=20,\ p=0.025$
\item $\nu=12,\ p=0.90$
\item $\nu=8,\ p=0.10$
\end{enumerate}
\end{multicols}

\vspace{5pt}{\large\bfseries\color{azul}5. Distribución F de Fisher--Snedecor}\par\vspace{3pt}

Encuentre \(F_{\nu_1,\nu_2,p}\); \(\nu_1\) corresponde al numerador y
\(\nu_2\) al denominador. Use las tablas de cola derecha \(0.05\) y
\(0.025\).

\begin{multicols}{2}
\begin{enumerate}[label=\alph*)]
\item $\nu_1=5,\nu_2=10,\ p=0.95$
\item $\nu_1=8,\nu_2=12,\ p=0.975$
\item $\nu_1=4,\nu_2=20,\ p=0.95$
\item $\nu_1=10,\nu_2=15,\ p=0.975$
\end{enumerate}
\end{multicols}

\newpage

\begin{center}
{\Large\bfseries Taller: aplicación de valores críticos}\\[-1pt]
{\small Cara B. Introducción numérica a los intervalos de confianza}
\end{center}

\textbf{Instrucciones.} Resuelva los cuatro problemas. En cada uno: (1)
identifique los datos y el nivel de confianza; (2) busque los
percentiles necesarios en la tabla; (3) sustituya en la fórmula
indicada; y (4) calcule ambos extremos del intervalo. No se requiere
demostrar las fórmulas.

\vspace{5pt}{\large\bfseries\color{azul}6. Intervalo para una media con desviación conocida (normal)}\par\vspace{3pt}

Una máquina llena envases con un contenido aproximadamente normal. En
una muestra aleatoria de \(n=64\) envases se obtuvo una media de
\(\bar x=502\) mL. Se conoce que la desviación estándar poblacional es
\(\sigma=16\) mL. Construya un \textbf{intervalo de confianza del 95\%}
para el contenido medio poblacional.

\textbf{Procedimiento:} consulte \(z_{0.975}\); calcule el error
estándar \(\sigma/\sqrt n\); obtenga el margen de error
\(E=z_{0.975}\sigma/\sqrt n\) y los límites \(\bar x-E\) y \(\bar x+E\).

\[IC_{95\%}(\mu)=\left[\bar x-z_{0.975}\frac{\sigma}{\sqrt n},\;\bar x+z_{0.975}\frac{\sigma}{\sqrt n}\right].\]

\vspace{5pt}{\large\bfseries\color{azul}7. Intervalo para una media con desviación desconocida (t)}\par\vspace{3pt}

Un laboratorio mide la resistencia a la tracción de \(16\) piezas
fabricadas bajo las mismas condiciones. La muestra presenta
\(\bar x=84\) MPa y desviación estándar muestral \(s=8\) MPa. Suponga
que la resistencia poblacional sigue una distribución normal. Construya
un \textbf{intervalo de confianza del 95\%} para la resistencia media.

\textbf{Procedimiento:} calcule \(\nu=n-1\); busque \(t_{\nu,0.975}\);
calcule \(s/\sqrt n\); determine el margen de error y los límites del
intervalo.

\[IC_{95\%}(\mu)=\left[\bar x-t_{\nu,0.975}\frac{s}{\sqrt n},\;\bar x+t_{\nu,0.975}\frac{s}{\sqrt n}\right].\]

\vspace{5pt}{\large\bfseries\color{azul}8. Intervalo para una varianza poblacional ($\chi^2$)}\par\vspace{3pt}

En un proceso de mecanizado se registraron los diámetros de \(n=21\)
piezas seleccionadas aleatoriamente. La varianza muestral fue
\(s^2=0.36\) mm\(^2\). Suponga normalidad e independencia. Construya un
\textbf{intervalo de confianza del 95\%} para la varianza poblacional
\(\sigma^2\).

\textbf{Procedimiento:} calcule \(\nu=n-1\); consulte
\(\chi^2_{\nu,0.025}\) y \(\chi^2_{\nu,0.975}\); calcule \((n-1)s^2\);
divida este numerador por cada valor crítico, respetando el orden de los
límites.

\[IC_{95\%}(\sigma^2)=\left[\frac{(n-1)s^2}{\chi^2_{\nu,0.975}},\;\frac{(n-1)s^2}{\chi^2_{\nu,0.025}}\right].\]

\vspace{5pt}{\large\bfseries\color{azul}9. Intervalo para una razón de varianzas (F)}\par\vspace{3pt}

Dos líneas independientes de producción elaboran componentes cuyas
mediciones se suponen normales. En la línea A, \(n_A=11\) y
\(s_A^2=25\); en la línea B, \(n_B=16\) y \(s_B^2=16\) (unidades
cuadradas). Construya un \textbf{intervalo de confianza del 95\%} para
\(\sigma_A^2/\sigma_B^2\).

\textbf{Procedimiento:} calcule \(\nu_1=n_A-1\), \(\nu_2=n_B-1\) y
\(R=s_A^2/s_B^2\); consulte \(F_{\nu_1,\nu_2,0.975}\) y
\(F_{\nu_2,\nu_1,0.975}\); use el recíproco para la cola inferior y
calcule los dos límites.

\[IC_{95\%}\!\left(\frac{\sigma_A^2}{\sigma_B^2}\right)=\left[\frac{R}{F_{\nu_1,\nu_2,0.975}},\;R\,F_{\nu_2,\nu_1,0.975}\right].\]

\end{document}
